\begin{abstract}
We estimate the leakage current of CMOS logic gates and of a 3-bit carry-lookahead adder (CLA) under process, voltage and temperature (PVT) variation with machine-learning models instead of SPICE, for the PTM 22\,nm high-performance technology at a nominal supply of 1\,V.
Four transistor-level models (NMOS and PMOS, ON and OFF) predict the four terminal currents of a single transistor from 10,000 SPICE samples per case; each is chosen from four model types (Gaussian process, hybrid analytical\,+\,MLP, residual MLP, physics-informed network) on a validation set of an 80:10:10 split.
Using only these models and the recorded SPICE node voltages, the gate leakage is assembled with Kirchhoff's current law on the supply node.
Over 3,600 gate simulations (five gates, 18 input cases, 200 corners) the mean accuracy is 98.87\,\% (median 99.18\,\%); the plain sum of the 33 gate leakages predicts the adder leakage over 12,800 simulations (64 vectors, 200 corners) with a mean accuracy of 99.00\,\% (median 99.17\,\%).
The 99\,\% target is therefore met for the circuit but not for the gates, and the circuit result benefits from a cancellation: summing isolated gate leakages alone caps the adder accuracy at 99.16\,\% even with perfect gate models.
The remaining gate error is concentrated in OFF transistors operated with the full supply across them, a corner of the input space that uniform sampling covers with less than 0.5\,\% of the training rows.
We also report the negative results that shaped the method: a silent \texttt{altermod} bug in the data generation, a plain current sum that fails by orders of magnitude, and a drain--source symmetry that the device model does not have.
% src: results/partB/gate_results.csv, results/partC/cla3_results.csv, results/analysis/report_numbers/rail_point_row_counts.csv; docs/REPORT_NOTES.md sections 7, 8, 9.1
\end{abstract}
\begin{IEEEkeywords}
leakage estimation, PVT variation, machine learning, BSIM4, logic gates, carry-lookahead adder
\end{IEEEkeywords}